% Zone „Das kennst du schon“ – aus bank/ebenen/zone.jsonl (zusammenbau v0.1)
\einheitenkopf[zone][Kennst du schon]{Das kennst du schon}
\setcounter{aufgabe}{0}

%% TODO Titel: Ich-kann-Satz fehlt in der Bank (Platzhalter: Kettenname) – Nr. 1
%% TODO Anweisung: ein Satz über der ersten Teilaufgabe fehlt in der Bank – Nr. 1
\begin{aufgabe}{Punkte im räumlichen Koordinatensystem lesen und eintragen, Koordinatenebenen und Achsen benennen, ein Schrägbild lesen (welche Kante liegt in welcher Koordinatenebene, welche Koordinate ist auf einer Seitenfläche konstant)}
\begin{teile}
\teil Lies die Koordinaten des Punktes $P$ ab.

\begin{ksys3}[xyz]\rpunkt{3}{2}{4}{P}\end{ksys3}
$P($ \leerfeld | \leerfeld | \leerfeld $)$
\teil Ein Quader mit Kanten parallel zu den Achsen hat die gegenüberliegenden Ecken $O(0 | 0 | 0)$ und $G(5 | 3 | 2{,}5)$. Welche Koordinate ist auf der Deckfläche konstant, mit welchem Wert?
\end{teile}
\end{aufgabe}

%% TODO Titel: Ich-kann-Satz fehlt in der Bank (Platzhalter: Kettenname) – Nr. 2
%% TODO Anweisung: ein Satz über der ersten Teilaufgabe fehlt in der Bank – Nr. 2
\begin{aufgabe}{Verbindungsvektor zweier Punkte bilden (Spitze minus Fuß), Vektoren addieren und mit einer Zahl vervielfachen, Ortsvektor und Punkt unterscheiden}
\begin{teile}
\teil $A(1 | 2 | 0)$, $B(4 | 3 | 5)$: $\overrightarrow{AB}$? $\overrightarrow{AB} = ($ \leerfeld | \leerfeld | \leerfeld $)$
\teil $2 \cdot (1 | -3 | 0{,}5) + (-1 | 4 | 2)$?
\end{teile}
\end{aufgabe}

%% TODO Titel: Ich-kann-Satz fehlt in der Bank (Platzhalter: Kettenname) – Nr. 3
%% TODO Anweisung: ein Satz über der ersten Teilaufgabe fehlt in der Bank – Nr. 3
\begin{aufgabe}{Kollinearität zweier Vektoren prüfen (ist der eine ein Vielfaches des anderen?)}
\begin{teile}
\teil Ist $(2 | 4 | 6)$ ein Vielfaches von $(1 | 2 | 3)$? Wenn ja: Faktor?
\teil Ist $(1{,}5 | 3 | -0{,}5)$ ein Vielfaches von $(3 | 6 | -1)$? Wenn ja: Faktor?
\end{teile}
\end{aufgabe}

%% TODO Titel: Ich-kann-Satz fehlt in der Bank (Platzhalter: Kettenname) – Nr. 4
%% TODO Anweisung: ein Satz über der ersten Teilaufgabe fehlt in der Bank – Nr. 4
\begin{aufgabe}{Skalarprodukt zweier Vektoren in Koordinaten berechnen und Orthogonalität als „Skalarprodukt null“ erkennen}
\begin{teile}
\teil $(1 | 2 | 3) \cdot (4 | 0 | 1)$?
\teil $(-3 | 2 | 0{,}5) \cdot (2 | 4 | -6)$?
\end{teile}
\end{aufgabe}

%% TODO Titel: Ich-kann-Satz fehlt in der Bank (Platzhalter: Kettenname) – Nr. 5
%% TODO Anweisung: ein Satz über der ersten Teilaufgabe fehlt in der Bank – Nr. 5
\begin{aufgabe}{Geradengleichung in Parameterform lesen (Stütz- und Richtungsvektor) und die Punktprobe an einer Geraden}
\begin{teile}
\teil Die Gerade $g$ hat die Gleichung $\vec x = (1 | 0 | 4) + t \cdot (2 | 1 | -1)$ mit $t \in \mathbb{R}$. Stützvektor, Richtungsvektor?
\teil Liegt $P(-1{,}5 | 3 | 0{,}5)$ auf der Geraden $\vec x = (0 | 2 | 1) + t \cdot (1{,}5 | -1 | 0{,}5)$, $t \in \mathbb{R}$?
\end{teile}
\end{aufgabe}

%% TODO Titel: Ich-kann-Satz fehlt in der Bank (Platzhalter: Kettenname) – Nr. 6
%% TODO Anweisung: ein Satz über der ersten Teilaufgabe fehlt in der Bank – Nr. 6
\begin{aufgabe}{Lineare Gleichungssysteme mit zwei und drei Unbekannten lösen, auch unterbestimmte (eine Unbekannte frei wählen, Vielfache als gleichwertige Lösungen erkennen)}
\begin{gleichungsraster}[2]
\gl{\text{$x + y = 5$ und $x - y = 1$}} &
\gl{\text{$x + 2y - z = 1$, $y + z = 4$, $z = 3$}} \\
\end{gleichungsraster}
\end{aufgabe}

%% TODO Titel: Ich-kann-Satz fehlt in der Bank (Platzhalter: Kettenname) – Nr. 7
%% TODO Anweisung: ein Satz über der ersten Teilaufgabe fehlt in der Bank – Nr. 7
\begin{aufgabe}{Eine lineare Gleichung als Punktmenge lesen}
\begin{teile}
\teil Liegt $(2 | 7)$ auf der Geraden $y = 3x + 1$?
\teil Die Gerade $2x + 3y = c$ ($c$ eine Zahl) geht durch $(0{,}5 | 3)$. $c$? $c = $ \leerfeld
\end{teile}
\end{aufgabe}

%% TODO Titel: Ich-kann-Satz fehlt in der Bank (Platzhalter: Kettenname) – Nr. 8
%% TODO Anweisung: ein Satz über der ersten Teilaufgabe fehlt in der Bank – Nr. 8
\begin{aufgabe}{Verbindungsvektor zweier Punkte bilden (Spitze minus Fuß), Vektoren addieren und mit einer Zahl vervielfachen, Ortsvektor und Punkt unterscheiden – Fehler finden}
\begin{teile}
\teil Lina rechnet für $A(2 | 5 | 1)$ und $B(3 | 1 | 4)$: $\overrightarrow{AB} = (2-3 | 5-1 | 1-4) = (-1 | 4 | -3)$. Finde den Fehler und rechne richtig.
\end{teile}
\end{aufgabe}

%% TODO Titel: Ich-kann-Satz fehlt in der Bank (Platzhalter: Kettenname) – Nr. 9
%% TODO Anweisung: ein Satz über der ersten Teilaufgabe fehlt in der Bank – Nr. 9
\begin{aufgabe}{Verbindungsvektor zweier Punkte bilden (Spitze minus Fuß), Vektoren addieren und mit einer Zahl vervielfachen, Ortsvektor und Punkt unterscheiden – selbst rechnen}
\begin{teile}
\teil $A(4 | 1 | 3)$, $B(1 | 2 | 7)$: $\overrightarrow{AB}$? $\overrightarrow{AB} = ($ \leerfeld | \leerfeld | \leerfeld $)$
\end{teile}
\end{aufgabe}
